\begin{theorem}[Schur 通道签名对分拆 Gibbs 律的有限层完全可逆性]\label{thm:conclusion-schur-gibbs-signature-complete-invertibility}
\leanverified{paper\_conclusion\_schur\_gibbs\_signature\_complete\_invertibility}
固定整数 \(q\ge 2\) 与分辨率 \(m\)。记
\[
p_{m,q}(\nu):=P^{\mathrm{part}}_{m,q}(\nu)
\qquad(\nu\vdash q),
\]
以及
\[
s_{m,q}(\lambda):=
\frac{T_{q,\lambda}(m)}{f^\lambda\,T_{q,(q)}(m)}
\qquad(\lambda\vdash q).
\]
则映射
\[
p_{m,q}=(p_{m,q}(\nu))_{\nu\vdash q}
\longmapsto
s_{m,q}=(s_{m,q}(\lambda))_{\lambda\vdash q}
\]
是线性同构。更精确地，
\[
s_{m,q}(\lambda)=\sum_{\nu\vdash q}p_{m,q}(\nu)\chi^\lambda(\nu),
\]
并且存在显式反演公式
\[
p_{m,q}(\mu)=\frac1{z_\mu}\sum_{\lambda\vdash q}\chi^\lambda(\mu)\,s_{m,q}(\lambda)
\qquad(\mu\vdash q).
\]
\end{theorem}
\begin{proof}
由正文中分拆 Gibbs 律与 Schur 通道序参量的定义，
\[
m_\lambda(m):=\mathbb E_{m,q}\chi^\lambda(\sigma)
=
\sum_{\nu\vdash q}p_{m,q}(\nu)\chi^\lambda(\nu).
\]
另一方面，
\[
\frac{T_{q,\lambda}(m)}{T_{q,(q)}(m)}=f^\lambda m_\lambda(m),
\]
故
\[
s_{m,q}(\lambda)=m_\lambda(m)
=
\sum_{\nu\vdash q}p_{m,q}(\nu)\chi^\lambda(\nu).
\]

再用对称群字符的列正交关系
\[
\sum_{\lambda\vdash q}\chi^\lambda(\mu)\chi^\lambda(\nu)=z_\mu\,\mathbf 1_{\mu=\nu},
\]
便有
\[
\sum_{\lambda\vdash q}\chi^\lambda(\mu)s_{m,q}(\lambda)
=
\sum_{\nu\vdash q}p_{m,q}(\nu)
\sum_{\lambda\vdash q}\chi^\lambda(\mu)\chi^\lambda(\nu)
=
z_\mu\,p_{m,q}(\mu),
\]
即得所示反演公式。由于正反公式都显式存在，该变换是线性同构。证毕。
\end{proof}

\begin{theorem}[字符签名的 Parseval 等距与相图单纯形]\label{thm:conclusion-schur-signature-parseval-simplex}
在分拆函数空间上定义加权 Hilbert 结构
\[
\langle a,b\rangle_z:=\sum_{\nu\vdash q}z_\nu\,a(\nu)\overline{b(\nu)},
\]
并定义变换
\[
\mathcal F_q:a\longmapsto \widehat a,
\qquad
\widehat a(\lambda):=\sum_{\nu\vdash q}a(\nu)\chi^\lambda(\nu).
\]
则有严格等距律
\[
\sum_{\lambda\vdash q}\widehat a(\lambda)\overline{\widehat b(\lambda)}
=
\sum_{\nu\vdash q}z_\nu\,a(\nu)\overline{b(\nu)}.
\]
特别地，对任意两个分拆概率向量 \(p,p'\) 及其 Schur 签名 \(s,s'\)，有
\[
\sum_{\lambda\vdash q}|s(\lambda)-s'(\lambda)|^2
=
\sum_{\nu\vdash q}z_\nu\,|p(\nu)-p'(\nu)|^2.
\]

再记
\[
v_\nu:=(\chi^\lambda(\nu))_{\lambda\vdash q}\in \CC^{\mathcal P(q)}.
\]
则所有可能的签名恰构成字符列向量的仿射单纯形
\[
\Sigma_q:=\operatorname{conv}\{v_\nu:\nu\vdash q\},
\]
并且任意签名
\[
s=\sum_{\nu\vdash q}p(\nu)v_\nu
\]
中的 \(p(\nu)\) 正是点 \(s\) 在单纯形 \(\Sigma_q\) 中的重心坐标。
\end{theorem}
\begin{proof}
由列正交关系，
\[
\sum_{\lambda\vdash q}\widehat a(\lambda)\overline{\widehat b(\lambda)}
=
\sum_{\mu,\nu\vdash q}
a(\mu)\overline{b(\nu)}
\sum_{\lambda\vdash q}\chi^\lambda(\mu)\chi^\lambda(\nu)
=
\sum_{\nu\vdash q}z_\nu\,a(\nu)\overline{b(\nu)}.
\]
取 \(a=b\) 即得等距公式；再令 \(a=p-p'\)，便得对概率向量的 Parseval 恒等式。

另一方面，由定理 \ref{thm:conclusion-schur-gibbs-signature-complete-invertibility}，
\[
s(\lambda)=\sum_{\nu\vdash q}p(\nu)\chi^\lambda(\nu),
\]
亦即
\[
s=\sum_{\nu\vdash q}p(\nu)v_\nu.
\]
由于 \(\mathcal F_q\) 可逆，列向量族 \(\{v_\nu\}_{\nu\vdash q}\) 线性无关，从而仿射无关；因此其凸包是一个仿射单纯形。重心坐标的唯一性正是上述可逆性的几何表达。证毕。
\end{proof}

\begin{theorem}[压力扇形在字符单纯形中的整胞塌缩]\label{thm:conclusion-pressure-fan-collapses-to-character-simplex-cells}
对每个压力向量 \(p=(P(1),\dots,P(q))\)，定义极大分拆集
\[
R_{\max}(p):=\{\nu\vdash q:\ P(\nu)=P_{\max}\}.
\]
并令
\[
\mathfrak S_q(p):=\lim_{m\to\infty}s_{m,q},
\]
其中极限按正文中的 pressure fan 定理存在。对任意非空集合 \(R\subset \mathcal P(q)\)，若胞腔
\[
\mathcal C_R:=\{p:\ R_{\max}(p)=R\}
\]
非空，则 \(\mathfrak S_q\) 在整块 \(\mathcal C_R\) 上为常值，并且
\[
\mathfrak S_q(p)=\sum_{\nu\in R}\omega_R(\nu)\,v_\nu,
\qquad
\omega_R(\nu):=
\frac{K(\nu)/z_\nu}{\sum_{\mu\in R}K(\mu)/z_\mu}.
\]
特别地：
\begin{enumerate}
\item 若 \(R=\{\nu^\ast\}\) 为单点，则
\[
\mathfrak S_q(p)=v_{\nu^\ast};
\]
\item 若 \(|R|\ge 2\)，则 \(\mathfrak S_q(p)\) 落在面
\[
\operatorname{conv}\{v_\nu:\nu\in R\}
\]
的相对内部。
\end{enumerate}
换言之，压力向量在 Schur 极限层面只留下“哪一组分拆达到最大值”这一组合标签，而不再留下胞腔内部的连续参数。
\end{theorem}
\begin{proof}
若 \(R_{\max}(p)=\{\nu^\ast\}\)，则正文中的纯冻结定理给出
\[
p_{m,q}\to \delta_{\nu^\ast}
\]
指数收敛。代入定理 \ref{thm:conclusion-schur-gibbs-signature-complete-invertibility} 的有限层表示，得
\[
s_{m,q}\to v_{\nu^\ast}.
\]

若 \(|R_{\max}(p)|\ge 2\)，则正文中的多极大分拆定理给出
\[
p_{m,q}(\nu)\to
\begin{cases}
\omega_R(\nu),& \nu\in R,\\
0,& \nu\notin R,
\end{cases}
\]
其中
\[
\omega_R(\nu)=\frac{K(\nu)/z_\nu}{\sum_{\mu\in R}K(\mu)/z_\mu}.
\]
再次代入定理 \ref{thm:conclusion-schur-gibbs-signature-complete-invertibility}，便得
\[
\mathfrak S_q(p)=\sum_{\nu\in R}\omega_R(\nu)v_\nu.
\]
由于每个 \(\omega_R(\nu)\) 都严格为正且总和为 \(1\)，故该点位于相应面之相对内部。右端只由 \(R\) 决定，因此 \(\mathfrak S_q\) 在整块 \(\mathcal C_R\) 上为常值。证毕。
\end{proof}

\begin{corollary}[唯一冻结相的指数可判定性]\label{cor:conclusion-pure-phase-nearest-vertex-exponential-decodability}
若 \(R_{\max}(p)=\{\nu^\ast\}\)，并记
\[
\Delta_q:=\min_{\mu\neq \nu}|v_\mu-v_\nu|_2>0,
\]
则存在常数 \(C_q,\delta_q>0\)，使得
\[
|s_{m,q}-v_{\nu^\ast}|_2\le C_q e^{-m\delta_q}.
\]
从而当 \(m\) 充分大时，\(\nu^\ast\) 是与 \(s_{m,q}\) 距离最近的唯一顶点。亦即，仅凭有限层 Schur 签名的最近顶点判别，即可最终精确恢复冻结主循环型。
\end{corollary}
\begin{proof}
由定理 \ref{thm:conclusion-pressure-fan-collapses-to-character-simplex-cells} 的纯相情形与正文中的指数收敛率，可得每个坐标都有
\[
s_{m,q}(\lambda)=\chi^\lambda(\nu^\ast)+O(e^{-m\delta^\ast}).
\]
由于分拆集合有限，故存在 \(C_q,\delta_q>0\) 使
\[
|s_{m,q}-v_{\nu^\ast}|_2\le C_q e^{-m\delta_q}.
\]
又 \(\{v_\nu\}_{\nu\vdash q}\) 为有限集，故最小点间距 \(\Delta_q\) 严正。若
\[
|s_{m,q}-v_{\nu^\ast}|_2<\frac{\Delta_q}{2},
\]
则三角不等式保证 \(v_{\nu^\ast}\) 为唯一最近顶点。证毕。
\end{proof}

\begin{theorem}[hook 通道生成函数的闭式因子化]\label{thm:conclusion-hook-channel-generating-polynomial-factorization}
对分拆
\[
\nu=1^{c_1}2^{c_2}\cdots q^{c_q}\vdash q,
\]
定义 hook 角色生成多项式
\[
H_\nu(t):=\sum_{k=0}^{q-1}\chi^{(q-k,1^k)}(\nu)\,t^k.
\]
则有严格恒等式
\[
H_\nu(t)=\frac{1}{1+t}\prod_{r=1}^{q}\bigl(1-(-t)^r\bigr)^{c_r}.
\]
并且映射
\[
\nu\longmapsto H_\nu(t)
\]
是单射。更具体地，若写
\[
(1+t)H_\nu(t)=\prod_{d=1}^{q}\Phi_d(-t)^{\,b_d},
\]
其中 \(\Phi_d\) 为第 \(d\) 个 cyclotomic 多项式，则
\[
b_d=\sum_{r:\ d\mid r}c_r,
\qquad
c_r=\sum_{k\ge 1}\mu(k)\,b_{kr}.
\]
\end{theorem}
\begin{proof}
令 \(V=\CC^q\) 为 \(S_q\) 的置换表示，\(W=V/\CC\mathbf 1\) 为标准表示。经典外幂恒等式给出
\[
\sum_{k=0}^{q-1}\chi_{\wedge^k W}(\nu)\,t^k
=
\det(I+t\,\sigma|_W),
\]
其中 \(\sigma\) 为循环型 \(\nu\) 的代表元。又
\[
\wedge^k W\cong S^{(q-k,1^k)},
\]
故左端正是 \(H_\nu(t)\)。

再由 \(V\cong \CC\mathbf 1\oplus W\)，有
\[
\det(I+t\,\sigma|_V)=(1+t)\det(I+t\,\sigma|_W).
\]
而在置换表示上，一个 \(r\)-循环贡献的特征多项式为
\[
\det(I+t\,\sigma_r)=1-(-t)^r.
\]
于是
\[
\det(I+t\,\sigma|_V)=\prod_{r=1}^{q}\bigl(1-(-t)^r\bigr)^{c_r},
\]
两式相除便得
\[
H_\nu(t)=\frac{1}{1+t}\prod_{r=1}^{q}\bigl(1-(-t)^r\bigr)^{c_r}.
\]

再把右端作 cyclotomic 分解：
\[
1-(-t)^r=\prod_{d\mid r}\Phi_d(-t).
\]
故
\[
(1+t)H_\nu(t)=\prod_{d=1}^{q}\Phi_d(-t)^{\,b_d},
\qquad
b_d=\sum_{r:\ d\mid r}c_r.
\]
由 Möbius 反演立得
\[
c_r=\sum_{k\ge 1}\mu(k)\,b_{kr}.
\]
因此 \(H_\nu(t)\) 唯一决定全部 \(c_r\)，从而唯一决定分拆 \(\nu\)。证毕。
\end{proof}

\begin{theorem}[纯相极限中非平凡 hook 通道的无损恢复性]\label{thm:conclusion-pure-phase-hook-family-complete-recovery}
定义 hook 归一化签名
\[
h_{m,q}(k):=
\frac{T_{q,(q-k,1^k)}(m)}
{f^{(q-k,1^k)}\,T_{q,(q)}(m)}
\qquad(0\le k\le q-1),
\]
以及其生成多项式
\[
H_{m,q}(t):=\sum_{k=0}^{q-1}h_{m,q}(k)\,t^k.
\]
则有精确有限层恒等式
\[
H_{m,q}(t)=\sum_{\nu\vdash q}p_{m,q}(\nu)\,H_\nu(t).
\]

若进一步 \(R_{\max}(p)=\{\nu^\ast\}\)，则存在 \(\delta^\ast>0\) 使得
\[
H_{m,q}(t)=H_{\nu^\ast}(t)+O(e^{-m\delta^\ast})
\]
按系数逐项指数收敛。因而当 \(m\) 充分大时，仅凭
\[
\bigl(h_{m,q}(1),\dots,h_{m,q}(q-1)\bigr)
\]
这 \(q-1\) 个非平凡 hook 通道，便可唯一恢复 \(\nu^\ast\)。
\end{theorem}
\begin{proof}
由定理 \ref{thm:conclusion-schur-gibbs-signature-complete-invertibility}，
\[
h_{m,q}(k)=\sum_{\nu\vdash q}p_{m,q}(\nu)\chi^{(q-k,1^k)}(\nu).
\]
对 \(k\) 求和即得
\[
H_{m,q}(t)=\sum_{\nu\vdash q}p_{m,q}(\nu)H_\nu(t).
\]

若 \(R_{\max}(p)=\{\nu^\ast\}\)，则正文中的纯冻结结果给出
\[
p_{m,q}\to \delta_{\nu^\ast}
\]
指数收敛，故
\[
H_{m,q}(t)\to H_{\nu^\ast}(t)
\]
亦按系数指数收敛。

再由定理 \ref{thm:conclusion-hook-channel-generating-polynomial-factorization}，映射 \(\nu\mapsto H_\nu(t)\) 为单射；而分拆集合有限，故不同模板间存在正最小距离。于是当 \(m\) 充分大时，\(H_{m,q}(t)\) 的最近模板必唯一等于 \(H_{\nu^\ast}(t)\)。又因为
\[
h_{m,q}(0)=\frac{T_{q,(q)}(m)}{T_{q,(q)}(m)}=1
\]
恒等成立，故生成多项式已由 \(q-1\) 个非平凡系数 \((h_{m,q}(1),\dots,h_{m,q}(q-1))\) 唯一决定，从而 \(\nu^\ast\) 也被唯一恢复。证毕。
\end{proof}

\begin{conjecture}[hook-stratum 分离]\label{conj:conclusion-hook-stratum-separation}
对每个固定 \(q\)，设
\[
\mathcal R_q:=\{R_{\max}(p):\ p\in \RR^q\}
\]
为全部可实现的极大分拆集。对每个 \(R\in\mathcal R_q\)，定义其 hook 极限模板
\[
\mathcal H_R(t):=
\sum_{\nu\in R}
\frac{K(\nu)/z_\nu}{\sum_{\mu\in R}K(\mu)/z_\mu}\,
H_\nu(t).
\]
则应有
\[
R\neq R'
\quad\Longrightarrow\quad
\mathcal H_R(t)\neq \mathcal H_{R'}(t).
\]
等价地，尽管 hook 投影在整个分拆单纯形上并不保持全局单射，但在真实出现的压力扇形胞腔上，它仍应保持层标签的完全可分。
\end{conjecture}

\noindent
以上结果表明，Schur 通道并非冻结相的附属观测，而是分拆 Gibbs 律的精确 Fourier 坐标：有限层上，分拆概率向量与 Schur 签名严格可逆；极限层上，整个 pressure fan 在字符单纯形中整胞塌缩为顶点与面心；而在纯相极限中，整套不可约谱又进一步坍缩到非平凡 hook 子族而仍保持主循环型的无损可识别性。于是，分拆 Gibbs 几何、Schur 断层坐标与 hook 最小通道族三者在结论层被压接为同一条刚性链。

\endinput
